\chapter{Resultados Parciais}
\label{cap:resultados}

Este capítulo apresenta os resultados já consolidados do estudo. O trabalho está organizado em etapas incrementais: a reprodução do baseline epidemiológico, a construção do grafo heterogêneo temporal, o treinamento da GNN temporal e a avaliação comparativa entre a GNN e o baseline SIR nos dois regimes descritos na metodologia. O objetivo é triplo: demonstrar que o pipeline reproduz fielmente o baseline de referência, evidenciar que o grafo construído possui estrutura suficiente para justificar a abordagem neural, e comparar os dois modelos tanto na reconstrução retroativa das trajetórias (Modo~1) quanto na previsão de curto horizonte fora da amostra (Modo~2). Adianta-se que os modelos são complementares: o SIR, por ajustar dois parâmetros a cada curva completa, é mais preciso na reconstrução \emph{in-sample}, ao passo que a GNN, um preditor global e congelado, supera o SIR de forma consistente e estatisticamente significativa na tarefa de previsão, que é o seu uso pretendido.

%=====================================================

\section{Reprodução do baseline SIR}

A primeira etapa teve como meta reproduzir o ajuste SIR do trabalho de referência \citep{oliveira2025brasnam} sobre os charts brasileiros, como pré-requisito metodológico para qualquer comparação posterior. O ajuste foi conduzido sobre o subconjunto \emph{viral$\cap$hit} (músicas presentes simultaneamente no Top 200 e no Viral 50), totalizando 1.981 músicas no pipeline local, valor próximo às 1.977 reportadas no trabalho original. Para cada música, as séries de viralidade e sucesso foram ajustadas de forma independente, estimando os parâmetros $\beta$, $\gamma$ e a razão $R_0 = \beta/\gamma$.

A Tabela~\ref{tab:phase0-r4} compara as métricas obtidas com os alvos do trabalho de referência. Todos os critérios mínimos pré-acordados foram atingidos: o RMSE de viralidade e de sucesso ficou dentro da margem de 10\% dos valores de referência, o teste de Mann-Whitney manteve-se na mesma ordem de grandeza, e o ajuste convergiu para 100\% das séries.

\begin{table}[htbp]
\centering
\caption{Critério de reprodução do baseline SIR.}
\label{tab:phase0-r4}
\begin{tabular}{p{4.5cm}p{2.5cm}p{2.5cm}p{2.0cm}}
\hline
\textbf{Métrica} & \textbf{Alvo (referência)} & \textbf{Obtido} & \textbf{Status} \\
\hline
RMSE viralidade & $\approx 0{,}028$ & $0{,}0289$ & Atingido \\
RMSE sucesso & $\approx 0{,}052$ & $0{,}0471$ & Atingido \\
Mann-Whitney $p$ & ordem $10^{-60}$ & $1{,}61 \times 10^{-39}$ & Atingido \\
Subset viral$\cap$hit & $1.977$ & $1.981$ & Atingido \\
Convergência do ajuste & $100\%$ & $100\%$ & Atingido \\
\hline
\end{tabular}
\end{table}

% FIGURA: réplica do boxplot da Fig. 3 do paper de referência (distribuição de RMSE por chart, SIR). (adiada)

A reprodução fiel cumpre uma função metodológica central: sem ela, a comparação posterior entre GNN e SIR não seria defensável, pois não haveria garantia de que o pipeline de dados e de avaliação está alinhado ao da literatura. Além de validar o pipeline, a análise das séries trouxe duas observações que orientam as próximas etapas. A primeira é que músicas com permanência longa nos charts (mais de cem dias ativos, cerca de 37\% do Top 200) exibem padrão multi-onda que eleva o RMSE médio, ainda que a mediana permaneça dentro da tolerância; esse comportamento justifica conceitualmente as extensões wave-based da literatura \citep{oliveira2025asonam}, mesmo que não tenham sido reimplementadas aqui. A segunda é a necessidade de investigar, antes da submissão, eventuais diferenças entre média e mediana na forma como o trabalho de referência reporta o erro.

%=====================================================

\section{Grafo heterogêneo temporal}

A segunda etapa construiu o grafo heterogêneo temporal descrito no Capítulo~\ref{cap:metodologia}. O grafo final contém 6.526 nós de música, 1.701 nós de artista e 530 nós de gênero. As contagens de música e gênero ficaram dentro das tolerâncias estabelecidas, e o número de artistas coincidiu exatamente com o esperado. O universo de música é maior que o subconjunto \emph{viral$\cap$hit} por decisão de projeto: incluir músicas presentes apenas no Top 200 enriquece o contexto estrutural disponível ao modelo.

A Tabela~\ref{tab:phase1-edges} resume as contagens por tipo de aresta. A relação \texttt{cotrajectory} domina o grafo, com mais de 660 mil arestas, refletindo o grande número de pares de músicas que compartilham trajetória nos charts. As demais relações (autoria, associação a gênero e coocorrência de gêneros) contribuem com volumes menores, porém semanticamente distintos.

\begin{table}[htbp]
\centering
\caption{Contagens por tipo de aresta no grafo heterogêneo.}
\label{tab:phase1-edges}
\begin{tabular}{p{7.0cm}r}
\hline
\textbf{Tipo de aresta} & \textbf{Contagem} \\
\hline
\texttt{artist} $\to$ \texttt{performs} $\to$ \texttt{music} & 9.274 \\
\texttt{artist} $\to$ \texttt{has\_genre} $\to$ \texttt{genre} & 3.344 \\
\texttt{music} $\to$ \texttt{cotrajectory} $\to$ \texttt{music} & 664.577 \\
\texttt{genre} $\to$ \texttt{cooccurs} $\to$ \texttt{genre} & 9.866 \\
\hline
\end{tabular}
\end{table}

A análise topológica indica que o grafo possui estrutura rica, e não conexões aleatórias. O coeficiente de \textit{clustering} médio é elevado tanto em \texttt{cotrajectory} (0,5126) quanto em \texttt{cooccurs} (0,6040), sugerindo agrupamentos coesos de músicas e de gêneros. A detecção de comunidades por Louvain sobre o subgrafo de gêneros identificou agrupamentos coerentes, com as cinco maiores comunidades reunindo 136, 119, 96, 60 e 43 gêneros, e apenas cinco comunidades isoladas de pequeno porte. As distribuições de grau são fortemente assimétricas, característica esperada em redes culturais, com poucos nós de altíssimo grau e uma longa cauda de nós pouco conectados.

% FIGURA: 4 painéis de distribuição de grau (um por tipo de aresta), escala log. (adiada)
% FIGURA: top-5 comunidades Louvain coloridas no subgrafo de gêneros. (adiada)

A construção foi acompanhada de um conjunto de validações automatizadas (nove critérios), que verificam desde a contagem de nós e a ausência de arestas órfãs até propriedades temporais e de integração com a arquitetura neural. Entre elas destacam-se a garantia de que todas as músicas do subconjunto \emph{viral$\cap$hit} estão presentes no grafo, a verificação de que o atributo \texttt{first\_seen\_week} permanece no intervalo válido em todas as arestas, e a confirmação de que o mecanismo de \textit{snapshots} temporais é monotônico, isto é, o grafo observável até uma semana está contido no grafo observável em semanas posteriores, propriedade essencial para evitar vazamento temporal. Um \emph{smoke-test} confirmou ainda que a arquitetura HeteroGraphSAGE percorre o grafo e produz embeddings de música com a dimensão esperada, sem valores inválidos. Todos os critérios foram satisfeitos.

%=====================================================

\section{Avaliação dupla: GNN temporal \emph{versus} SIR}
\label{sec:avaliacao-dupla}

A terceira etapa treinou a GNN heterogênea temporal (arquitetura HeteroGraphSAGE seguida de uma unidade recorrente GRU, janela temporal de doze semanas, dimensão latente 128 e três camadas de convolução) e a avaliou contra o baseline SIR sobre exatamente o mesmo eixo semanal e o mesmo subconjunto de músicas, condição indispensável para uma comparação justa. O modelo selecionado atingiu erro quadrático médio de validação de $7{,}5 \times 10^{-4}$. A avaliação segue dois regimes distintos e complementares. As pequenas variações no número de músicas entre etapas (1.981 no ajuste SIR, 1.955 no Modo~1 e 1.933 no Modo~2) refletem exigências mínimas específicas de cada regime (por exemplo, uma janela causal completa de $W$ semanas antes da origem de previsão), de modo que uma música é excluída apenas da etapa em que a exigência correspondente não é satisfeita.

\subsection{Modo 1: reconstrução retroativa}

No Modo~1, cada modelo reconstrói a trajetória semanal inteira de uma música a partir de uma semente inicial. Trata-se de uma tarefa \emph{in-sample} para o SIR, cujos parâmetros $\beta$ e $\gamma$ são ajustados individualmente à curva completa de cada música; a GNN, ao contrário, é um modelo global e congelado que simula a trajetória às cegas. A Tabela~\ref{tab:phase3-mode1} mostra que, nesse cenário, o SIR obtém menor RMSE em ambos os regimes, tanto na métrica de espectro completo quanto na restrita às semanas em que a música figura efetivamente no chart (\emph{on-chart}).

\begin{table}[htbp]
\centering
\caption{Modo 1: RMSE médio da reconstrução retroativa (1.955 músicas por regime).}
\label{tab:phase3-mode1}
\begin{tabular}{lrrrr}
\hline
\textbf{Regime} & \textbf{GNN} & \textbf{SIR} & \textbf{GNN (on-chart)} & \textbf{SIR (on-chart)} \\
\hline
Viral 50 & $0{,}0266$ & $0{,}0141$ & $0{,}1915$ & $0{,}1756$ \\
Top 200  & $0{,}0632$ & $0{,}0343$ & $0{,}1637$ & $0{,}1262$ \\
\hline
\end{tabular}
\end{table}

Essa vantagem do SIR é esperada e não contradiz a hipótese do trabalho: comparar um modelo global congelado, que prevê a partir de uma semente de doze semanas, contra um ajuste de dois parâmetros livres feito sobre a curva inteira favorece estruturalmente o segundo. O ponto relevante é que a diferença se estreita fortemente quando se olham apenas as semanas \emph{on-chart}, onde a dinâmica de fato importa: no regime Viral 50 os erros ficam próximos ($0{,}1915$ contra $0{,}1756$). O critério pré-registrado de superioridade da GNN na reconstrução (denominado C4 no protocolo) não foi atingido e é, portanto, documentado honestamente como uma limitação do regime \emph{in-sample}, e não como falha do modelo. A mesma leitura se aplica ao subgrupo de \emph{hits} de longa permanência (mais de noventa dias em chart, 806 músicas), no qual o SIR mantém vantagem (RMSE médio $0{,}0605$ contra $0{,}0998$ da GNN): são justamente as trajetórias multi-onda que o ajuste curva-a-curva melhor captura.

\subsection{Modo 2: previsão de curto horizonte}

O Modo~2 mede a tarefa de fato pretendida para a GNN: previsão causal, a $k$ semanas à frente, usando somente informação disponível até a origem da previsão. Aqui o SIR é reajustado de forma causal: $\beta$ e $\gamma$ são estimados por mínimos quadrados não lineares apenas sobre os pontos observados até a semana de origem, e a curva prevista resulta da integração da EDO por $k$ semanas a partir do último estado estimado, o que o coloca em pé de igualdade metodológica com a GNN. Incluiu-se ainda um baseline de persistência (repetir o último valor observado). A Tabela~\ref{tab:phase3-mode2} apresenta o RMSE por horizonte e regime.

\begin{table}[htbp]
\centering
\caption{Modo 2: RMSE da previsão causal a $k$ semanas (1.933 músicas pareadas por célula).}
\label{tab:phase3-mode2}
\begin{tabular}{llrrr}
\hline
\textbf{Regime} & \textbf{$k$} & \textbf{GNN} & \textbf{SIR} & \textbf{Persist.} \\
\hline
Viral 50 & 1 & $\mathbf{0{,}0261}$ & $0{,}0368$ & $0{,}0267$ \\
         & 2 & $\mathbf{0{,}0345}$ & $0{,}0604$ & $0{,}0375$ \\
         & 4 & $\mathbf{0{,}0444}$ & $0{,}0798$ & $0{,}0538$ \\
Top 200  & 1 & $0{,}0238$ & $0{,}0458$ & $\mathbf{0{,}0231}$ \\
         & 2 & $0{,}0370$ & $0{,}0524$ & $\mathbf{0{,}0360}$ \\
         & 4 & $\mathbf{0{,}0513}$ & $0{,}0636$ & $0{,}0520$ \\
\hline
\end{tabular}
\end{table}

A GNN supera o SIR em \emph{todos} os seis cenários (dois regimes $\times$ três horizontes). A vantagem é estatisticamente significativa: o teste de Wilcoxon pareado por música rejeita a hipótese nula em todas as células, com valores-$p$ entre $6{,}5 \times 10^{-4}$ e $5{,}4 \times 10^{-32}$, e a GNN apresenta menor erro em 64\% a 77\% das músicas, conforme o horizonte. A margem sobre o SIR também cresce com o horizonte no regime Viral 50 (de $29\%$ em $k=1$ para $44\%$ em $k=4$), indicando que a estrutura do grafo é especialmente útil quando a extrapolação epidemiológica se degrada. O baseline de persistência é forte nos horizontes curtos (comportamento esperado, pois a popularidade varia pouco de uma semana para a seguinte) e, no regime Top~200, supera marginalmente a GNN tanto em $k=1$ ($0{,}0231$ contra $0{,}0238$) quanto em $k=2$ ($0{,}0360$ contra $0{,}0370$); as diferenças são pequenas, mas documenta-se honestamente que o critério mínimo pré-registrado de superioridade sobre a persistência não é satisfeito em todas as células, no mesmo espírito da leitura dada ao critério C4 no Modo~1. No regime Viral~50, a GNN supera a persistência em todos os horizontes avaliados; no Top~200, apenas no horizonte mais longo ($k=4$: $0{,}0513$ contra $0{,}0520$). Frente ao SIR, por outro lado, a GNN é consistentemente superior nos seis cenários avaliados. Esse resultado, a GNN como o melhor preditor frente ao baseline epidemiológico, constitui a contribuição central do estudo e satisfaz o critério de previsão pré-registrado (C6) na comparação com o SIR; frente à persistência, o critério é atingido no Viral~50, mas apenas parcialmente no Top~200, onde a GNN só a supera em $k=4$.

%=====================================================

\section{Discussão parcial}

Os resultados consolidados sustentam três conclusões parciais. Primeiro, o baseline epidemiológico foi reproduzido dentro das margens da literatura, validando o pipeline de dados e avaliação. Segundo, a topologia do grafo heterogêneo (\textit{clustering} elevado, comunidades de gênero coerentes e distribuições de grau realistas) sugere que existe sinal estrutural a ser explorado por uma GNN. Terceiro, e mais importante, a avaliação dupla mostra que GNN e SIR são complementares, cada um superior na tarefa para a qual é estruturalmente adequado.

A leitura conjunta dos dois modos é o ponto central. Na reconstrução retroativa (Modo~1), o SIR leva vantagem, mas por um motivo metodológico transparente: ele ajusta dois parâmetros livres à curva completa de cada música, enquanto a GNN simula às cegas a partir de uma semente curta. Já na previsão causal de curto horizonte (Modo~2), que é o uso pretendido de um preditor e a pergunta de pesquisa efetivamente relevante, a GNN supera o SIR em todos os horizontes e regimes, com significância estatística forte. Em outras palavras, um preditor neural treinado sobre esse grafo supera o modelo epidemiológico puramente temporal, com vantagem que cresce à medida que o horizonte de previsão aumenta. Atribuir esse ganho especificamente à estrutura relacional (coautoria, gênero e cotrajetória), e não apenas à maior flexibilidade de um modelo neural frente a um ajuste paramétrico de dois graus de liberdade, é, neste estágio, hipótese e não achado estabelecido.

Resta como trabalho em andamento testar diretamente essa hipótese, por meio da análise de interpretabilidade da GNN (ablação por tipo de aresta, importância de grupos de atributos e extração de análogos de $\beta$, $\gamma$ e $R_0$ a partir das trajetórias previstas) e de um baseline neural sem grafo (por exemplo, uma GRU sobre a série de cada música isolada), que isola o ganho atribuível à flexibilidade do modelo daquele atribuível à estrutura relacional. Os resultados completos serão incorporados na versão final da dissertação.
